% ==============================================================================
% notes.tex -- Notes Added in Proof (pp. 505-510)
% ==============================================================================

\chapter*{Notes Added in Proof}
\addcontentsline{toc}{chapter}{Notes Added in Proof}
\markboth{\MakeUppercase{Notes Added in Proof}}{\MakeUppercase{Notes Added in Proof}}

\subsection*{Notes on Chapter 2}
\addcontentsline{toc}{section}{Notes on Chapter 2}

Let $E/S$ be an elliptic curve, $N \ge 1$ and $M \ge 1$ two integers, and $P,Q \in E[NM](S)$. Then we have
\[
(e_{NM}(P,Q))^M = e_{NM}(MP,Q) = E_N(MP,MQ) ,
\]
the last equality by applying (\ref{eq:2.8.4.1}) to the situation
\begin{gather*}
E = E_0 = E_1 = E_2,\quad \pi_1 = N,\quad \pi_2 = M, \\
P_0 = MP \in \operatorname{Ker}(\pi_1),\quad P_2 = Q \in \operatorname{Ker}(\pi_1^t \circ \pi_2^t) .
\end{gather*}
This compatibility occurs implicitly in (\ref{thm:2.9.1})'s description of $\mathcal{B}_N$.
Taking $M = 2$ and $P = Q$, we find
\[
1 = (e_{2N}(P,P))^2 = e_N(2P,2P), \text{ for } P \in E[2N](S) ,
\]
the first equality by (\ref{item:2.8.3}). Because $[2]: E \to E$ is f.p.p.f.\ surjective, any point $R$ in $E[N](S)$ is locally f.p.p.f.\ of the form $2P$ for $P \in E[2N](S)$, whence we have
\[
e_N(R,R) = 1 \text{ for } R \in E[N](S) .
\]

\subsection*{Notes on Chapter 4}
\addcontentsline{toc}{section}{Notes on Chapter 4}

Contrary to what is implicitly asserted in Section~\ref{sec:4.12}, the property ``being of given dimension'' is not in general local for the etale topology, or even for the Zariski topology (e.g., the spectrum of a discrete valuation ring, or the disjoint union of two schemes of different dimensions). It is rather the property ``non-empty, and every non-empty Zariski open set is $n$-dimensional'' which is local for the etale topology, and it is in this stronger sense that the phrase ``$\mathcal{P}$ is $n$-dimensional'' is to be interpreted, for $\mathcal{P}$ a relatively representable moduli problem on $(\mathrm{Ell})$, or on $(\mathrm{Ell}/k)$ for $k$ a field.

For a scheme which is locally of finite type over $\mathbb{Z}$ or over a field, ``$n$-dimensional'' in this stronger sense is equivalent to the property ``equidimensional, and $n$-dimensional'', i.e., ``non-empty, and every irreducible component is $n$-dimensional''; this last property, for such schemes, is equivalent to the property ``non-empty, and the local ring at every closed point is $n$-dimensional'' (cf.~EGA~IV, 10.6.1~(i), 10.6.2, 10.7.1, as well as 5.2.5(ii)). The point here is that for a scheme $X$ which is locally of finite type over $\mathbb{Z}$ (resp., over a specified field $k$), a point $x \in X$ is closed in $X$ if and only its residue field is a finite field (resp., a finite extension of $k$); thus for any open neighborhood $U$ of $x$, $x \in U \subset X$, the point $x$ is closed in $U$ if and only if it [is] closed in $X$. [It is this property which ``fails'' for $X$ the spectrum of a discrete valuation ring.]

If $f: X \to Y$ is a finite flat surjective morphism of locally noetherian schemes, the property ``non-empty, and the local ring at every closed point is $n$-dimensional'' holds for $X$ if and only if it holds for $Y$. [Because $f$ is finite, $x \in X$ is closed in $X$ if and only if $y = f(x)$ is closed in $Y$. Because $f$ is finite surjective, the inverse image of $\operatorname{Spec}(\mathcal{O}_{Y,y})$ is the disjoint union of the $\operatorname{Spec}(\widehat{\mathcal{O}}_{X,x_i})$ over the finite non-empty set of points $x_i \in X$ with $f(x_i) = y$. Because $f$ is finite and flat, each $\operatorname{Spec}(\widehat{\mathcal{O}}_{X,x_i})$ is individually finite and flat over $\operatorname{Spec}(\widehat{\mathcal{O}}_{Y,y})$, and so necessarily surjective as well, because its open image contains the closed point. By (EGA~IV, 5.4.1~(ii)), we conclude that $\dim(\widehat{\mathcal{O}}_{X,x}) = \dim(\widehat{\mathcal{O}}_{Y,y})$ for $x \in X$, $y = f(x)$. Applying this to closed points, we have the result.]

Similarly, for $f: X \to Y$ as above, the property ``equidimensional and $n$-dimensional'' holds for $X$ if and only if it holds for $Y$. [Apply the equality $\dim(\widehat{\mathcal{O}}_{X,x}) = \dim(\widehat{\mathcal{O}}_{Y,f(x)})$, proven above, to points where this dimension is zero to conclude that $x \in X$ is a maximal point of $X$ (i.e., a generic point of an irreducible component of $X$) if and only if $y = f(x)$ is a maximal point of $Y$. As $f$ is proper as well, this means that if $X_\alpha$ is an irreducible component of $X$, then $f(X_\alpha)$ is an irreducible component, say $Y_\beta$, of $Y$, and every irreducible component of $Y$ is of this form. Because the induced map $X_\alpha^{\mathrm{red}} \to Y_\beta$ is finite and surjective, we have $\dim(X_\alpha) = \dim(Y_\beta)$, again by EGA~IV, 5.4.1~(ii).]

Applying either of these equivalences, we see that for $\mathcal{P}$ a relatively representable moduli problem on $(\mathrm{Ell})$ (respectively, on $(\mathrm{Ell}/k)$ for $k$ a field) which is surjective finite and flat over $(\mathrm{Ell})$ (resp.~over $(\mathrm{Ell}/k)$), then $\mathcal{P}$ is automatically ``two-dimensional'' (resp., ``one-dimensional'') in the strong sense. [For if $E/S$ is a modular family, then $S$ is locally a smooth curve over $\mathbb{Z}$ (resp., over $k$), and $\mathcal{P}_{E/S} \to S$ is finite flat surjective by hypothesis.] For such a $\mathcal{P}$, if $\mathcal{P}'$ is a second relatively representable moduli problem on $(\mathrm{Ell})$ (resp.~on $(\mathrm{Ell}/k)$) which is finite over $(\mathrm{Ell})$ (resp.~over $(\mathrm{Ell}/k)$), and if $\mathcal{P} \to \mathcal{P}'$ is a finite flat surjective morphism of relatively representable moduli problems, then $\mathcal{P}'$ is itself finite flat surjective over $(\mathrm{Ell})$ (resp.~over $(\mathrm{Ell}/k)$), by a standard sorite on faithfully flat morphisms (cf.~Bourbaki, \emph{Comm.~Alg.}~I, \S3.4, Prop.~7). Therefore $\mathcal{P}'$ is itself ``two-dimensional'' (resp., ``one-dimensional'') in the strong sense. If in addition $\mathcal{P}$ is regular, then $\mathcal{P}'$ is also regular, being surjective and flat ``under'' $\mathcal{P}$ (cf.~EGA~IV, 6.5.2). This argument occurs in proving (\ref{thm:6.6.1}) and (\ref{thm:7.5.1}).

If $\mathcal{P}$ and $\mathcal{P}'$ are two relatively representable moduli problems on $(\mathrm{Ell})$ (resp.~on $(\mathrm{Ell}/k)$), both of which are regular and ``$n$-dimensional'' in the strong sense, then any finite morphism of moduli problems $\mathcal{P} \to \mathcal{P}'$ is automatically finite and flat. This observation is used repeatedly throughout the book, always accompanied by the catch-phrase ``a finite map between regular schemes of the same dimension is automatically flat.'' This catch-phrase is trivially false as stated unless the schemes in question are equidimensional as well (think of the disjoint union of a line and a point mapping to a line). The correct statement is ``a finite morphism $f: X \to Y$ between locally noetherian schemes both of which are regular and equidimensional of dimension $n$ is automatically flat.'' [Proof: Because $Y$ is regular and locally noetherian, it is the disjoint union of its irreducible components $Y_\alpha$, each of which is open and closed in $Y$. Looking at the situation over each $Y_\alpha$ with $f^{-1}(Y_\alpha)$ non-empty, we reduce to the case $Y$ irreducible. Because $X$ is regular and locally noetherian, it is the disjoint union of its irreducible components $X_\beta$, each open and closed in $X$, so we may examine the situation on each $X_\beta$ individually. Thus we may and will assume both $X$ and $Y$ irreducible, locally noetherian, regular $n$-dimensional. Because $f: X \to Y$ is finite, we have $\dim(X) \le \dim f(X) \le \dim(Y)$, by EGA~IV, 5.4.1~(i). Therefore $f$ is surjective. Combining EGA~IV, 5.6.4 with 5.6.5.3, we see that, because $f$ is finite, we have $\dim(\mathcal{O}_{X,x}) = \dim(\mathcal{O}_{Y,f(x)})$. The result now follows, by [EGA~IV, 15.4.2, $e') \implies b)$]. Indeed, the same chain of reasoning and references gives the same conclusion if one replaces ``finite'' by ``quasi-finite'' and ``$X$ regular'' by ``$X$ Cohen-Macaulay'' in the hypotheses.] The reader will check that wherever the catch-phrase ``a finite map between regular schemes of the same dimension is automatically flat'' appears in the text, it is applied either to schemes which are in fact equidimensional, or to moduli problems which have already been proven to be ``$n$-dimensional'' in the strong sense (cf.~(\ref{thm:1.12.7}), (\ref{thm:2.3.1}), (\ref{thm:5.5.2}), (\ref{cor:5.5.3}), (\ref{prop:7.5.2}), (\ref{thm:7.9.6}), (\ref{thm:10.12.2}), (\ref{thm:12.6.1}), (\ref{thm:12.8.2}), (\ref{thm:13.2.2})).

\subsection*{Notes on Chapters 8 and 10}
\addcontentsline{toc}{section}{Notes on Chapters 8 and 10}

The smoothness of coarse moduli schemes outside supersingular points in ``bad'' characteristics (cf.~(\ref{thm:10.9.7}), (\ref{thm:10.10.3})~(5), (\ref{cor:10.10.6})~(2)) is more easily proven by viewing these schemes as quotients of suitable fine moduli schemes (cf.~(\ref{item:8.1.1})), and applying the following quite general theorem.

\begin{theorem*}
Let $S$ be a regular noetherian scheme, $X \to S$ a smooth affine curve over $S$, and $G$ a finite group operating $S$-linearly on $X$. Then the quotient $X/G$ is a smooth curve over $S$.
\end{theorem*}

\begin{proof}
It suffices to check smoothness at points which are closed in their fibers. Localizing, completing, and ``algebraically closing the residue field'' (cf.~EGA~$0_{\mathrm{III}}$, 10.3.1 and 10.3.1~(5)) at a point $s \in S$, we reduce to the case $S = \operatorname{Spec}(A)$, with $A$ a complete noetherian regular local ring with algebraically closed residue field. Then the complete local ring of $X$ at a closed point $x$ of its special fiber $X_s$ is $A$-isomorphic to $A[[T]]$, the isotropy group $G_x \subset G$ of $x$ acts $A$-linearly on $A[[T]]$, and the complete local ring of $X/G$ at the image of $x$ is the ring of invariants of $G_x$ acting on $A[[T]]$.

The theorem therefore results from the following proposition (compare (\ref{prop:7.5.2})), applied to the image of $G_x$ in $\operatorname{Aut}_A(A[[T]])$.
\end{proof}

\begin{proposition*}
Let $A$ be a complete noetherian local domain, $T$ an indeterminate, and $G \subset \operatorname{Aut}_A(A[[T]])$ a finite group of $A$-linear automorphisms. Then
\[
(A[[T]])^G = A[[N_G(T)]] .
\]
\end{proposition*}

\begin{proof}
We have obvious inclusions
\[
A[[T]] \supset (A[[T]])^G \supset A[[N_G(T)]] .
\]
Just as in (\ref{prop:7.5.2}), we see that $A[[T]]$ is spanned over $A[[N_G(T)]]$ by the elements $1, T, \dots, T^{\#G-1}$. Therefore these elements certainly span $A[[T]]$ over $(A[[T]])^G$. By Galois theory, these elements must be linearly independent over $(A[[T]])^G$, so a fortiori linearly independent over $A[[N_G(T)]]$. Therefore $A[[T]]$ is a free module with basis $1, T, \dots, T^{\#G-1}$ over both $(A[[T]])^G$ and over $A[[N_G(T)]]$, whence the inclusion
\[
(A[[T]])^G \supset A[[N_G(T)]]
\]
must be an equality.
\end{proof}

Related to the discussion of coarse base change in Chapter~8 (cf.~(\ref{prop:8.1.6})), we have the following theorem.

\begin{theorem*}
Let $S$ be a regular noetherian scheme, $X \to S$ a smooth affine curve over $S$, and $G$ a finite group operating $S$-linearly on $X$. Suppose that for every geometric point $s$ of $S$, the action of $G$ on the fiber $X_{\bar{s}}$ is free on an open dense set of $X_{\bar{s}}$. Then for any regular noetherian $S'$ and any morphism $S' \to S$, the canonical map
\[
(X \times_S S')/G \to (X/G) \times_S S' \text{ is an isomorphism.}
\]
\end{theorem*}

\begin{proof}
The question is Zariski local on $S'$, which we may suppose connected, and hence irreducible. Source and target are smooth curves over $S'$ by the previous theorem, so both are regular noetherian schemes, equidimensional of the same dimension. By general properties of quotients, the map in question is finite, surjective (and radicial), so we may conclude that it is finite, flat, and surjective. To see that it is of degree one over $(X/G) \times_S S'$, it suffices to check over the dense open set where the action is free.
\end{proof}
